\documentclass{beamer}

\mode<presentation>
{
  \usetheme{Singapore}      % or try Darmstadt, Madrid, Warsaw, ...
  \usecolortheme{crane} % or try albatross, beaver, crane, ...
  \usefonttheme{serif}  % or try serif, structurebold, ...
  \setbeamertemplate{navigation symbols}{}
  \setbeamertemplate{caption}[numbered]
} 

\usepackage[autostyle=true]{csquotes}
\usepackage{xeCJK}
\usepackage[T1]{fontenc}
\usepackage{xcolor}
\usepackage{setspace}
\setmainfont[AutoFakeSlant=0.2]{Charis SIL}
\setCJKmainfont[AutoFakeSlant=0.2]{Noto Serif CJK HK}
\setCJKmonofont[AutoFakeSlant=0.2]{Noto Serif CJK HK}
\usepackage{polyglossia}
\setmainlanguage{english}
\usepackage{csquotes}
\usepackage{array}
\usepackage{makecell}

\title[Pearl River Sinitic]{Sinitic Languages of the Pearl River Region \\ --- Unit 2.2}
\subtitle{13ᵗʰ EACL Summer School, Universität Stuttgart}
\author{Hilário de Sousa 蘇沙}
\institute{EHESS–CRLAO \\ {\tiny(membre associé)}}
\date{Tuesday 28 July 2026 \\ 丙午六月十五 \\ {\footnotesize \texttt{hilario@bambooradical.com}}}

\begin{document}

\begin{frame}
  \titlepage
\end{frame}

\section{Middle Chinese}

\begin{frame}[fragile]{Middle Chinese}
\enquote{Early Middle Chinese}:\\
切韻 cit3 wan5 / qièyùn (rhyme dictionary; 601)\\
\smallskip
{\footnotesize 隋 Suí (581-618) ended \(\approx \)300 years of disunity of China in 589\\
\quad 陸法言 luk6 faat3 jin4 / lù fǎyán gathered eight scholars from northern / southern background, and decided on the \enquote{best} pronunciation of characters\\
organised by finals > tones > initials
}
\medskip
反切 faan2 cit3 / fǎnqiè, e.g.\\
深 sam1: 式 sik1 針 zam1 切\\
\medskip

陳澧 can4 lai5 (1810—1882) 《切韻考》「反切繫聯法」\\
\smallskip
深 sam1: 式 sik1 針 zam1 切\\
式 sik1: 賞 soeng2 職 zik1 切\\
賞 soeng2: 書 syu1 兩 loeng5 切 etc...\\
\smallskip
[深式賞書... share the same initial]\\

\end{frame}

\end{document}
